Dans une expérience d'effet photoélectrique, on éclaire successivement une cathode avec deux lumières monochromatiques et on mesure le potentiel d'arrêt :
\begin{center}
\begin{tabular}{cc}
Longueur d'onde & Potentiel d'arrêt \\ \hline
\qty{\v{lam1}}{nm} & \qty[round-mode=places,round-precision=3]{\fpeval{1240/\v{lam1} - \v{phi}}}{V} \\
\qty{\v{lam2}}{nm} & \qty[round-mode=places,round-precision=3]{\fpeval{1240/\v{lam2} - \v{phi}}}{V}
\end{tabular}
\end{center}

\begin{enumerate}[label=\alph*)]
    \item À partir de ces deux mesures seulement, déterminez la valeur expérimentale de la constante de Planck.
    \item Déterminez le travail d'extraction de la cathode, en électronvolts.
    \item Tracez l'allure du graphique du potentiel d'arrêt en fonction de la fréquence. Que représentent la pente et l'abscisse à l'origine ?
\end{enumerate}
